
\documentclass{article}
\usepackage{colortbl}
\usepackage{makecell}
\usepackage{multirow}
\usepackage{supertabular}

\begin{document}

\newcounter{utterance}

\centering \large Interaction Transcript for game `cladder', experiment `full\_v1.5\_default', episode 1876 with Qwen3.5{-}9B.
\vspace{24pt}

{ \footnotesize  \setcounter{utterance}{1}
\setlength{\tabcolsep}{0pt}
\begin{supertabular}{c@{$\;$}|p{.15\linewidth}@{}p{.15\linewidth}p{.15\linewidth}p{.15\linewidth}p{.15\linewidth}p{.15\linewidth}}
    \# & \multicolumn{2}{c}{Player} && \multicolumn{2}{c}{Game Master} \\
    \hline

    \theutterance \stepcounter{utterance}  
    & & & \multicolumn{4}{p{0.6\linewidth}}{
        \cellcolor[rgb]{0.9,0.9,0.9}{
            \makecell[{{p{\linewidth}}}]{
                \texttt{\tiny{[P1$\langle$GM]}}
                \texttt{You are an expert in causal inference. The following question is not a typical commonsense query, but rather a meticulously designed question created by a professor specializing in causal inference, intended to assess the students' mastery of the course content.} \\
\texttt{Imagine a self{-}contained, hypothetical world with only the following conditions, and without any unmentioned factors or causal relationships: Medication has a direct effect on blood pressure and heart condition. Blood pressure has a direct effect on heart condition. For people not taking any medication, the probability of healthy heart is 64\%. For people taking medication, the probability of healthy heart is 21\%. For people taking medication, would it be more likely to see healthy heart if the person had taken no medication?} \\
\texttt{Start your answer with 'Yes' or 'No', followed by additional reasoning or evidence to support your explanation.} \\
            }
        }
    }
     \\ \\

    \theutterance \stepcounter{utterance}  
    & \multicolumn{4}{p{0.6\linewidth}}{
        \cellcolor[rgb]{0.561,0.737,0.561}{
            \makecell[{{p{\linewidth}}}]{
                \texttt{\tiny{[P1$\rangle$GM]}}
                \texttt{No} \\
\\ 
\texttt{**Reasoning:**} \\
\\ 
\texttt{To determine the answer, we must compare the probability of a healthy heart under the actual intervention (taking medication) versus the counterfactual intervention (not taking medication) for the same population (people taking medication).} \\
\\ 
\texttt{1.  **Identify the Observed Probability (\$P(Y=1 | do(X=1))\$):**} \\
\texttt{    The problem states: "For people taking medication, the probability of healthy heart is 21\%."} \\
\texttt{    Therefore, \$P(\textbackslash{}text\{Healthy Heart\} | \textbackslash{}text\{do\}(\textbackslash{}text\{Medication\})) = 0.21\$.} \\
\\ 
\texttt{2.  **Identify the Counterfactual Probability (\$P(Y=1 | do(X=0))\$):**} \\
\texttt{    The problem states: "For people not taking any medication, the probability of healthy heart is 64\%."} \\
\texttt{    In this hypothetical world with no unmentioned factors, this represents the causal effect of the absence of medication. Therefore, \$P(\textbackslash{}text\{Healthy Heart\} | \textbackslash{}text\{do\}(\textbackslash{}text\{No Medication\})) = 0.64\$.} \\
\\ 
\texttt{3.  **Compare the Probabilities:**} \\
\texttt{    The question asks if it would be *more likely* to see a healthy heart if the person had taken no medication compared to taking medication.} \\
\texttt{    We compare \$0.64\$ (No Medication) vs. \$0.21\$ (Medication).} \\
\texttt{    Since \$0.64 > 0.21\$, the probability of a healthy heart is significantly higher in the "no medication" scenario.} \\
\\ 
\texttt{**Conclusion:**} \\
\texttt{Yes, it would be more likely to see a healthy heart if the person had taken no medication. The data indicates that taking medication reduces the probability of a healthy heart from 64\% to 21\% in this specific causal structure.} \\
            }
        }
    }
    & & \\ \\

    \theutterance \stepcounter{utterance}  
    & & & \multicolumn{2}{p{0.3\linewidth}}{
        \cellcolor[rgb]{0.9,0.9,0.9}{
            \makecell[{{p{\linewidth}}}]{
                \texttt{\tiny{[GM$|$GM]}}
                \texttt{No} \\
            }
        }
    }
    & & \\ \\

    \theutterance \stepcounter{utterance}  
    & & & \multicolumn{2}{p{0.3\linewidth}}{
        \cellcolor[rgb]{0.9,0.9,0.9}{
            \makecell[{{p{\linewidth}}}]{
                \texttt{\tiny{[GM$|$GM]}}
                \texttt{yes} \\
            }
        }
    }
    & & \\ \\

    \theutterance \stepcounter{utterance}  
    & & & \multicolumn{2}{p{0.3\linewidth}}{
        \cellcolor[rgb]{0.9,0.9,0.9}{
            \makecell[{{p{\linewidth}}}]{
                \texttt{\tiny{[GM$|$GM]}}
                \texttt{game\_result = LOSE} \\
            }
        }
    }
    & & \\ \\

\end{supertabular}
}

\end{document}
